% Source: https://artofproblemsolving.com/community/c3h2101465_geometry
% Stable identity: aops-topic-2101465-q01
Let $\theta=\angle XAB=\angle YAC$ and $\alpha=\angle BAC$, and let $P$ and $Q$ be the midpoints of $AB$ and $AC$.

\begin{asy}
usepackage("amsmath");
size(7cm);
defaultpen(fontsize(10pt));
pair A=(1.6,3.4), B=(0,0), C=(5.2,0);
real th=32;
pair ux=dir(degrees(B-A)-th), uy=dir(degrees(C-A)+th);
pair B0=A+dot(B-A,ux)*ux, C0=A+dot(C-A,uy)*uy;
pair M=(B+C)/2, P=(A+B)/2, Q=(A+C)/2;
draw(A--B--C--cycle,linewidth(0.9));
draw(A--(A+1.45*(B0-A)),gray); draw(A--(A+1.4*(C0-A)),gray);
label("$X$",A+1.45*(B0-A),W); label("$Y$",A+1.4*(C0-A),E);
draw(B--B0,dashed); draw(C--C0,dashed);
draw(P--B0,deepgreen); draw(P--M,deepgreen); draw(Q--M,deepgreen); draw(Q--C0,deepgreen);
draw(M--B0,red+linewidth(1)); draw(M--C0,red+linewidth(1));
draw(arc(A,0.8,degrees(B0-A),degrees(B-A)),blue); draw(arc(A,0.8,degrees(C-A),degrees(C0-A)),blue);
label("$\theta$",A+0.95*dir((degrees(B0-A)+degrees(B-A))/2),blue+fontsize(8pt));
label("$\theta$",A+0.95*dir((degrees(C-A)+degrees(C0-A))/2),blue+fontsize(8pt));
dot(A); dot(B); dot(C); dot(B0); dot(C0); dot(M); dot(P); dot(Q);
label("$A$",A,N); label("$B$",B,SW); label("$C$",C,SE); label("$M$",M,S);
label("$B_0$",B0,dir(100)); label("$C_0$",C0,dir(80)); label("$P$",P,dir(150)); label("$Q$",Q,dir(30));
label("\mbox{}",(min(currentpicture,user=true)+max(currentpicture,user=true))/2,black); // keeps the text after the figure black when LaTeX typesets the labels inline
\end{asy}

\emph{Lengths.} By the midline theorem, $PM\parallel AC$ with $PM=\dfrac{AC}{2}$, and $QM\parallel AB$ with $QM=\dfrac{AB}{2}$. Since $\theta<90^\circ$, the foot $B_0$ lies on the ray $AX$, and triangle $AB_0B$ has a right angle at $B_0$; its median to the hypotenuse is half the hypotenuse, so $PB_0=PA=PB=\dfrac{AB}{2}$. Similarly $QC_0=QA=QC=\dfrac{AC}{2}$. Hence
\[PB_0=QM\qquad\text{and}\qquad PM=QC_0 .\]

\emph{Angles.} Triangle $APB_0$ is isosceles with $PA=PB_0$, so $\angle PB_0A=\angle PAB_0=\theta$, and the exterior angle at $P$ is $\angle BPB_0=2\theta$. Since $PM\parallel AC$, $\angle BPM=\angle BAC=\alpha$. The ray $AX$ lies on the other side of line $AB$ from $C$ (otherwise angle $XAB$ would contain interior points of the triangle), so $B_0$ and $M$ are on opposite sides of line $AB$. Hence the angle from $PM$ through $PB$ to $PB_0$ is $\alpha+2\theta$, and $\angle MPB_0$ equals $\alpha+2\theta$ if this is at most $180^\circ$, and $360^\circ-(\alpha+2\theta)$ otherwise.

In exactly the same way, $\angle CQC_0=2\theta$, $\angle CQM=\alpha$ (as $QM\parallel AB$), $C_0$ and $M$ are on opposite sides of line $AC$, and $\angle MQC_0$ is $\alpha+2\theta$ or $360^\circ-(\alpha+2\theta)$ by the same rule. So $\angle B_0PM=\angle MQC_0$.

\emph{Conclusion.} Triangles $B_0PM$ and $MQC_0$ have $PB_0=QM$, $PM=QC_0$ and equal included angles at $P$ and $Q$, so they are congruent (SAS). Therefore $MB_0=C_0M$. $\blacksquare$
